\section{Experimental Setup}
\label{sec:experimental.setup}
\subsubsection{Computing System} 
The experiments have been carried out on an Intel Xeon Gold 6248R processor with a Cascade micro-architecture. The processor consists of $24$ cores operating at a base frequency of $3$ GHz for non-AVX512 instructions.
%Supporting AVX-512 instruction sets, the processor can perform two fused multiply-add instructions with 8 double-precision floating-point numbers every clock cycle.
With $24$ cores and a peak AVX-512 boost frequency of $2.5$ GHz, the processor achieves a theoretical data throughput of ca. $1.92$ double precision TFlops/s
We measured a peak performance of $1.78$ double precision Tflops/s using the likwid performance tool \tq.
%todo paper 

The source code has been compiled with \tf{GCC} \tf{v10.2} using the highest optimization level \tf{-Ofast} and \tf{-march=native}, \tf{-pthread} and \tf{-fopenmp}. 
Loops within for the eighth case have been parallelized using \tf{GCC}'s \tf{OpenMP} \tf{v4.5} implementation.
We have used the \tf{GEMV} and \tf{GEMM} implementation of the 2024.0 Intel oneAPI Math Kernel Library (oneMKL) and its own threading library \tf{mkl\_intel\_thread} together with the threading runtime library \tf{libiomp5}.
The benchmark results of each function are the average of 10 runs.

\subsubsection{Tensor Shapes} 
We have used asymmetrically-shaped and symmetrically-shaped tensors in order to cover many possible use cases. 
The dimension tuples of both shape types are organized within two three-dimensional arrays with which tensors are initialized.
The dimension array for the first shape type contains $720 = 9\times 8 \times 10$ dimension tuples where the row number is the tensor order ranging from $2$ to $10$. 
For each tensor order $8$ tensor instances with increasing tensor size is generated.
%todo: explain specialities of the first set: first dimension is ??, one dimension is greater than ?? and other dimensions are equal to 2.
The second set consists of $336 = 6\times8\times 7$ dimensions tuples where the tensor order ranges from $2$ to $7$ and has $8$ dimension tuples for each order.
Each tensor dimension within the second set is $2^{12}$, $2^{8}$, $2^{6}$, $2^5$, $2^4$ and $2^3$.
A detailed explanation of the tensor shape setup is given in \cite{bassoy:2019:ttv, bassoy:2018:fast}.

%todo: explain why you chose certain visiualization.
%todo: do I need contour plots or can I take 
\begin{comment}
The number of runtime values equals to the number of elements of the tensor sets.
Instead of providing all runtime data points, the arithmetic mean over the set of tensor sizes is given.
Hence, the performance values are provided in three dimensional performance plot
that vary with the tensor order and contraction mode resulting in a three dimensional performance plot.
Hence, countour view (performance maps) are provided.
\end{comment}

%\tit{Setup 1} performs runtime measurements with \tit{asymmetrically}-\tit{shaped} tensors while \tit{Setup 2} performs runtime measurements with \tit{symmetrically}-\tit{shaped} tensors.

\begin{comment}
Their dimension tuples are organized in $10$ two-dimensional arrays $\mbN_q$ with $9$ rows and $32$ columns where the dimension tuple $\mbn_{r,c}$ of length $r+1$ denotes an element $\mbN_q(r,c)$ of $\mbN_q$ with $1 \leq q \leq 10$.
The dimension $\mbn_{r,c}(i)$ of $\mbN_q$ is $1024$ if $i = 1$, $c \cdot 2^{15-r}$ if $i = \min(r+1,q)$ and $2$ for any other index $i$ with $1 < q \leq 10$.
The dimension $\mbn_{r,c}(i)$ of $\mbN_1$ is given by $c \cdot 2^{15-r}$ if $i=1$, $1024$ if $i=2$ and $2$ for any other index $i$.
Dimension tuples of the same array column have the same number of tensor elements.
Please note that with increasing tensor order (and row-number), the contraction mode is halved and with increasing tensor size, the contraction mode is multiplied by the column number.
Such a setup enables an orthogonal test-set in terms of tensor elements ranging from $2^{25}$ to $2^{29}$ and tensor order ranging from $2$ to $10$.
\tit{Setup 2} performs runtime measurements with \tit{symmetrically}-\tit{shaped} tensors.
Their dimension tuples are organized in one two-dimensional array $\mbM$ with $6$ rows and $8$ columns where the dimension tuple $\mbm_{r,c}$ of length $r+1$ denotes an element $\mbM(r,c)$ of $\mbM$.
For $c=1$, the dimensions of $\mbm_{r,c}$ are given by $2^{12}$, $2^{8}$, $2^{6}$, $2^5$, $2^4$ and $2^3$ with descending row number $r$ from $6$ to $1$.
For $c>1$, the remaining dimensions are given by $\mbm_{r,c} = \mbm_{r,c} + k \cdot (c-1)$ where $k$ is $2^9$, $2^5$, $2^3$, $2^2$, $2$, $1$ with descending row number $r$ from $6$ to $1$.
In this setup, shape tuples of a column do not yield the same number of subtensor elements.
\end{comment}


\begin{comment}
\subsubsection{Performance Maps} 
%In this setup, a mode-$q$ tensor-vector multiplication is executed with the shape array $\mbN_q$.
% where the contraction dimension is equal to the largest dimension.
Measuring a single tensor-vector multiplication with the first setup produces $2880 = 9\times32\times 10$ runtime data points where the tensor order ranges from $2$ to $10$, with $32$ shapes for each order and $10$ contraction modes.
The second setup produces $336 = 6\times8\times 7$ data points with $6$ tensor orders ranging from $2$ to $7$, $8$ shapes for each order and $7$ contraction modes.
Similar to the findings in \cite{bassoy:2018:fast}, we have observed a performance loss for small dimensions of the mode with the highest priority.
% is small with a standard deviation where the performance values deviate from the mean between $5$\% and $20$\%.
The presented performance values are the arithmetic mean over the set of tensor sizes that vary with the tensor order and contraction mode resulting in a three dimensional performance plot.
A schematic countour view of the plots is given in Fig. \ref{fig:performance.map} which is divided into $5$ regions.
The cases $2$, $3$, $6$ and $7$ generate performance values within the regions $2$, $3$, $6$ and $7$ where only a single parallel \tf{GEMV} is executed, see Table \ref{tab:mapping}.
Please note that the contraction mode $q$ is set to the tensor order $p$ if $q>p$.
Performance values within region $8$ result from case $8$ which executes \tf{GEMV}'s with tensor slices in parallel.

%The contraction mode $q$ is set equal to the tensor order $p$ for all measurements if $q>p$.
%There are a lot of data runtime measurement data.
%A function is measured with 2880 and 480 different shape tuples for one data type.
\begin{figure*}[t]
	\centering
	\includegraphics[width=0.5\textwidth]{figures/performance_map_color}
	\caption{
		\footnotesize %
		Schematic contour view of the following average performance maps for the tensor-vector multiplication with tensors that are stored according to the first-order storage format.
		Each case $x$ in Table \ref{tab:mapping} affects a different region $x$ within the performance map.
		Performance values are the arithmetic mean over the set of tensor sizes with $32$ and $8$ elements in case of the first and second test setup, respectively.
		Contraction mode $q=p$ for $q>p$ where $p$ is the tensor order.
		\label{fig:performance.map}
	}
\end{figure*}
The following analysis considers four parallel versions \tf{SB-P1}, \tf{LB-P1}, \tf{SB-PN} and \tf{LB-PN}.
\tf{SB} (small-block) and \tf{LB} (large-block) denote parallel slice-vector multiplications where each thread recursively calls a single-threaded \tf{GEMV} with mode-$2$ and mode-$\mhq$ slices, respectively.
\tf{P1} uses the outer-most dimension $n_{p}$ for parallel execution whereas \tf{PN} applies loop fusion and considers all fusible dimensions for parallel execution.
%Figure Schematic contour view of the following average performance maps
%\subsubsection{Variations of Case $8$} 
%According to Table \ref{tab:mapping}
%executing \tf{GEMV}s with tensor slices in parallel is given by the upper triangular covering the $8$-th case. 
%Average performance values of the cases $2$, $3$, $6$ and $7$ are covered by the remaining parts of the map.
\end{comment}